\documentclass[a4paper]{article}

%%%% Packages %%%%
\usepackage{amsfonts}
\usepackage{amsmath}
\usepackage{amssymb}
\usepackage{amsthm}
\usepackage{bbm}
\usepackage{booktabs}
\usepackage{calc}
\usepackage{caption}
\usepackage{centernot}
\usepackage{circuitikz}
\usepackage{dsfont}
\usepackage{esint}
\usepackage{esvect}
\usepackage{float}
\usepackage{fullpage}
\usepackage{graphicx}
\usepackage{hhline}
\usepackage{hyperref}
\usepackage{ifthen}
\usepackage{lipsum}
\usepackage{marvosym}
\usepackage{mathrsfs}
\usepackage{mathtools}
\usepackage{microtype}
\usepackage{multicol}
\usepackage{parskip}
\usepackage{siunitx}
\usepackage{svg}
\usepackage{tikz}
\usepackage{tikz-cd}
\usepackage{titlesec}
\usepackage{titling}
\usepackage{wrapfig}
\usepackage{xcolor}
\usepackage[inline]{enumitem}
\usepackage[makeroom]{cancel}
\usepackage[most]{tcolorbox}
\usetikzlibrary{angles}
\usetikzlibrary{arrows.meta}
\usetikzlibrary{bending}
\usetikzlibrary{calc}
\usetikzlibrary{patterns}
\usetikzlibrary{decorations.pathmorphing}
\usetikzlibrary{decorations.markings}

\setsvg{inkscapeexe="C:/Program Files/Inkscape/bin/inkscape.exe"}

\hypersetup{
	pdfborder={0 0 0}
}
\makeatletter
\renewcommand{\@makefnmark}{\hbox{\textsuperscript{\color{red}\normalfont\@thefnmark}}}
\makeatother
%%%% Packages %%%%
%%%% Custom Commands %%%%

\DeclareMathOperator{\Div}{div}

% Define a problem counter that resets with the section
\newcounter{problem}[section]

% Redefine the problem numbering to include section and problem numbers
\renewcommand{\theproblem}{\arabic{problem}}

% Define the Problem command - BG and ENG versions
\newcommand{\Problem}{\refstepcounter{problem}\fbox{\textbf{Problem \theproblem:}} }
\newcommand{\ProblemBG}{\refstepcounter{problem}\fbox{\textbf{Задача \theproblem:}} }

% Solution Definition
\newcommand{\Solution}[1]{\textbf{Solution\ifthenelse{\equal{#1}{}}{}{ #1}:}}
\newcommand{\SolutionBG}[1]{\textbf{Решение\ifthenelse{\equal{#1}{}}{}{ #1}:}}

% This forces LaTeX to reserve vertical space as if every line had a fraction - use in \begin{aligned} .. \end{aligned}
\newcommand{\rowstrut}{\vphantom{\dfrac{1}{2}}}
%%%% Custom Commands %%%%
% Resistor
\tikzset{
	resistor/.style={
		draw=none,
		decorate,
		decoration={
			markings,
			mark=at position 0.5 with {
				% Connectors
				\draw[line] (-\pgfdecoratedcompleteddistance,0) -- (-0.75,0);
				\draw[line] (0.75,0) -- (\pgfdecoratedcompleteddistance,0);
				% Resistor
				\draw[line] (-0.75, 0.35) rectangle (0.75, -0.35);
				\node[above, transform shape] at (0, 0.35) {$R$};
			}
		}
	}
}

% Electrical battery (POSSIBLY FIX WITH RESISTOR CODE FOR OLYMPIAD 22)
\tikzset{
	battery/.style={
		draw=none,
		decorate,
		decoration={
			markings,
			mark=at position 0.5 with {
				% Connectors
				\draw[line] (-\pgfdecoratedcompleteddistance,0) -- (-4,0);
				\draw[line] (4,0) -- (\pgfdecoratedcompleteddistance,0);
				% Battery
				\draw[line] (-4, -20) -- (-4, 20);
				\draw[line] (4, -10) -- (4, 10);
				% Polarity
				\node[above] at (-4, -15) {$+$};
				\node[below] at (4, -15) {$-$};
			}
		}
	}
}

% Electrical switch
\tikzset{
	switch/.style={
		draw=none,
		decorate,
		decoration={
			markings,
			mark=at position 0.5 with {
				% Connectors
				\draw[line] (-\pgfdecoratedcompleteddistance,0) -- (-15,0);
				\draw[line] (15,0) -- (\pgfdecoratedcompleteddistance,0);
				% Switch
				\draw[line] (-12.5,0) circle (2.5);
				\draw[line] (12.5, 0) circle (2.5);
				\draw[line] (-10,0) -- (10,7);
			}
		}
	}
}

% Springs or inductors
\tikzset{
	spring/.style={
		decorate,
		decoration={
			coil,
			aspect=0.6,
			segment length=7pt,
			amplitude=6pt,
			pre length=5pt,
			post length=5pt
		},
		line width=1.5
	}
}

% A point for electrical terminals (radius = inner sep + line width / 2)
% USAGE: \node[terminal, xshift=-2.5pt] at (1, 1) {};
\tikzset{
	terminal/.style={
		circle,
		draw=black,
		line width=1.5pt,
		inner sep=1.75pt,
		append after command={
			\pgfextra{
				\edef\terminalname{\tikzlastnode}
				\draw[black, line width=1.5pt]
				([shift={(-4.5pt,-4.5pt)}]\terminalname.center) --
				([shift={(-1.94pt,-1.94pt)}]\terminalname.center);
				\draw[black, line width=1.5pt]
				([shift={(1.94pt,1.94pt)}]\terminalname.center) --
				([shift={(4.5pt,4.5pt)}]\terminalname.center);
			}
		}
	}
}

% Just a point for junctions
\tikzset{
	point/.style={
		circle,
		fill=black,
		inner sep=1.5pt
	}
}

% Standardized line
\tikzset{
	line/.style={
		line width=1.5pt,
	}
}


%%% THIS IS A HELPER GRID FOR TIKZ PICTURES %%%

%\begin{center}
%	\begin{tikzpicture}
%		
%		% Grid
%		\draw[help lines, step=1] (-2,-2) grid (12,12);
%		
%		% Axes
%		\draw[->, line width=1] (-2,0) -- (12.5,0) node[right] {$x$};
%		\draw[->, line width=1] (0,-2) -- (0,12.5) node[above] {$y$};
%		
%		% Axis labels
%		\foreach \x in {-2,-1,1,2,...,12}
%		\node[below] at (\x,0) {\small $\x$};
%		
%		\foreach \y in {-2,-1,1,2,...,12}
%		\node[left] at (0,\y) {\small $\y$};
%		
%		% Your drawing
%		\draw[spring] (0,0) -- (10,10);
%		
%	\end{tikzpicture}
%\end{center}



%%% ROTATING ARROW HEADS %%%

% To rotate {Stealth} in place use the flex' command:
%\draw[
%line,
%-{Stealth[length=7pt,flex'=0.9]}] 
%(175,186.5) arc (90:270:12.5);



%%% EXAMPLE HATCHING FOR SURFACES %%%
%\begin{scope}
%	\clip (350,100) rectangle (425,275);
%	
%	\foreach \x in {175,183,...,425} {
%		\draw[line width=0.75pt]
%		(\x,100) -- ++(175,+175);
%	}
%\end{scope}



%%% EXAMPLE ANGLE ARC $$$
% Angles
%\draw[line]
%(7,3)
%arc[start angle=-180, end angle=0, radius=2];
\input{patterns_kogan}

%\usepackage[T2A]{fontenc}
%\usepackage[utf8]{inputenc}
\usepackage[bulgarian]{babel}

\usepackage[
left=2.5cm,
right=2.5cm,
top=3cm,
bottom=2.5cm
]{geometry}

\setlength{\droptitle}{-7em}
\title{\Huge Б. Ю. Коган - 100 Задачи по Физика}
\author{}
\date{}


\begin{document}
	\maketitle
	
	\ProblemBG % Коган Задача 001
	\textbf{\underline{Равнопроменливо Движение $\# 1$}}
	\vspace{0.2em}
	\\
	Материална точка, намираща се в началния момент в състояние на покой, започва да се движи равноускорително. На края на първия километър скоростта ѝ е равна на $10\,\mathrm{m/s}$.
	\begin{description}[leftmargin=!, labelwidth=\widthof{$\bullet$}]
		\item[$\bullet$] Да се намери израз $\Delta v_{n, n+1}$ за промяната на скоростта между $n$-тия и $(n+1)$-вия километър.
		\item[$\bullet$] Ползвайки получения израз за $\Delta v_{n, n+1}$, намерете скоростите $v_1, v_2, \cdots, v_{10}$ съответно на $1$-вия, $2$-рия, $\cdots$, $10$-тия километър. Да се построи графика на скоростта спрямо пътя $(v[\mathrm{km/h}]\,\,\text{vs.}\,\,s[\mathrm{km}])$.
		\item[$\bullet$] В тази подточка ще ползваме числените стойности, които сме табулирали от предната част:
		\begin{description}[leftmargin=!, labelwidth=\widthof{$\rightarrow$}]
			\item[$\rightarrow$] Числените стойности за $s[\mathrm{km}]$ може да се приемат за стойностите на координатата $x[\mathrm{km}]$;
			\item[$\rightarrow$] Числените стойности за $v[\mathrm{km/h}]$ може да се приемат за стойностите на координатата $y[\mathrm{km}]$.
		\end{description}
		Получената графика $(y[\mathrm{km}]\,\,\text{vs.}\,\,x[\mathrm{km}])$ е стробоскопична снимка на тяло хвърлено от началото на координатната система под ъгъл в хипотетично нехомогенно гравитационно поле\footnote{Да се приеме като математически модел, а не физически постижимо поле: целта е експериментално упражнение.}:
		\begin{equation*}
			g(y) = \frac{\gamma}{y^3}
		\end{equation*} 
		където $\gamma$ е неизвестна константа. Знае се, че снимките се вземат през интервал от време $\Delta t = 20\,\mathrm{s}$.
		Да се намери константата $\gamma$. Определете скоростта и ъгъла на изстрелване на частицата.
	\end{description}
	
	\ProblemBG % Коган Задача 002
	\textbf{\underline{Равнопроменливо Движение $\# 2$}}
	\vspace{0.2em}
	\\
	Камък, хвърлен вертикално нагоре, се движи във вакуум в хомогенно гравитационно поле $g = 9{,}81\,\mathrm{m/s^2}$. За да достигне височина $h = 29{,}4\,\mathrm{m}$ за време $\Delta t_1 = 3\,\mathrm{s}$ е необходима начална скорост $v_1 = 24{,}5\,\mathrm{m/s}$, а за да направи това за $\Delta t_2 = 6\,\mathrm{s}$ $\rightarrow$ $v_2 = 34{,}3\,\mathrm{m/s}$. Защо за по-бързото достигане на една и съща височина се изисква по-малка начална скорост, а не по-голяма?
	
	\ProblemBG % Коган Задача 003
	\textbf{\underline{Равнопроменливо Движение $\# 3$}}
	\vspace{0.2em}
	\\
	Материална точка се движи равнозакъснително. В даден момент нейната скорост е $v = 10\,\mathrm{m/s}$. Да се намери при какво ускорение от този момент нататък частицата ще измине $s = 8\,\mathrm{m}$ за $\Delta t = 2\,\mathrm{s}$.
	
	\ProblemBG % Коган Задача 006
	\textbf{\underline{Относително Движение $\# 1$}}
	\vspace{0.2em}
	\\
	Две кръгли платформи, разположени една до друга и всяка с радиус $R = 2\,\mathrm{m}$, се въртят в противоположни посоки, както е показано на фигурата. Разстоянието между геометричните им центрове е равно на $d = 5\,\mathrm{m}$, а ъгловата скорост на всяка платформа е равна на $\omega = 1\,\mathrm{rad/s}$. На перифериите на двете платформи се намират наблюдатели $A_1$ и $A_2$, които заемат в дадения момент положенията, показани на чертежа. Ако всеки от наблюдателите е неподвижен спрямо частта от платформата, на която е стъпил, то с каква скорост се движи наблюдателят $A_2$ по отношение на наблюдателя $A_1$? Как се променя относителната скорост, ако и двамата наблюдатели във всеки един момент от време гледат в едно и също фиксирано направление, примерно по посока на $y$?
	\begin{figure}[H]
		\centering
		\includesvg[width=0.49\textwidth]{figures_svg/kogan_6}
	\end{figure}
	
	\ProblemBG % Коган Задача 005
	\textbf{\underline{Относително Движение $\# 2$}}
	\vspace{0.2em}
	\\
	Наблюдателят $A$ се движи с постоянна по модул скорост $v_A$. Наблюдателят $B$ се движи също с постоянна по модул скорост $v_B$, която е по-голяма от $v_A$ и във всеки момент от време има същата посока като нея. Може ли при тези условия наблюдателят $B$ да се окаже неподвижен спрямо наблюдателя $A$? Можете ли да илюстрирате с пример или контрапример?
	
	\ProblemBG % Коган Задача 008
	\textbf{\underline{Miscellaneous $\# 1$}}
	\vspace{0.2em}
	\\
	При движението си дъска, поставена на кръгли ролки, запазва хоризонтално положение и остава на една и съща височина. Може ли да се осъществи същото движение, ползвайки некръгли ролки?
	
	\ProblemBG % Коган Задача 009
	\textbf{\underline{Miscellaneous $\# 2$}}
	\vspace{0.2em}
	\\
	На везни е уравновесен сандък, напълнен със сачми. Ще останат ли везните в равновесие, ако се заменят сачмите с по-едри (направени от същия материал)? Приемете термодинамична граница:
	\begin{description}[leftmargin=!, labelwidth=\widthof{$\rightarrow$}]
		\item[$\rightarrow$] Размерите на сандъка са много по-големи от диаметъра на сачмите.
		\item[$\rightarrow$] Процентното увеличение на сачмите не е достатъчно голямо, за да наруши горното условие.
		\item[$\rightarrow$] Сачмите са многобройни, така че заеманият от тях обем не зависи от точното им разположение\footnote{При малък сандък и големи сачми може, в зависимост от начина на разположение, да поберем различен брой сачми. Работим в рамките на термодинамичната граница именно, за да пренебрегнем този ефект (packing problems).}.
	\end{description}
	
	\ProblemBG % Коган Задача 010
	\textbf{\underline{Статика $\# 1$}}
	\vspace{0.2em}
	\\
	Към дъска, лежаща на подпорите $A$ и $B$, са приложени силите $F_1$ и $F_2$, както е показано на фигурата. Ще се измени ли огъването на дъската, ако двете сили се заменят с равнодействуващата им сила $R$?
	\begin{figure}[H]
		\centering
		\includesvg[width=0.6\textwidth]{figures_svg/kogan_10}
	\end{figure}
	
	\ProblemBG % Коган Задача 014
	\textbf{\underline{Влак}}
	\vspace{0.2em}
	\\
	На фигурата е изобразен локомотив, движещ се надясно. В положението, в което е показана фигурата, мотовилката (коляно-мотовилковият механизъм) действува на колелото с някаква сила, насочена наляво и минаваща над точка $C$. Но тази точка е център, около който в дадения момент се върти колелото. Затова под действие на тази сила то трябва да се завърти по посока, обратна на часовниковата стрелка, т.е. да се завърти наляво. Защо локомотивът се движи надясно?
	\begin{figure}[H]
		\centering
		\includesvg[width=0.835\textwidth]{figures_svg/kogan_14}
	\end{figure}
	
	\newpage
\section*{Решения}
\setcounter{problem}{0}

\SolutionBG{1}
	
	\newpage
	\section*{Ползвани Задачи}
	\begin{description}
		\item[$\bullet$] Задача 001 $\rightarrow$ 001
		\item[$\bullet$] Задача 002 $\rightarrow$ 002
		\item[$\bullet$] Задача 003 $\rightarrow$ 003
		\item[$\bullet$] Задача 004 $\rightarrow$ 005
		\item[$\bullet$] Задача 005 $\rightarrow$ 006
		\item[$\bullet$] Задача 006 $\rightarrow$ 008
		\item[$\bullet$] Задача 007 $\rightarrow$ 009
		\item[$\bullet$] Задача 008 $\rightarrow$ 010
		\item[$\bullet$] Задача 009 $\rightarrow$ 014
	\end{description}
\end{document}